\documentclass[11pt]{article}
\usepackage[a4paper,margin=24mm]{geometry}
\usepackage{hyperref}
\usepackage{stepdiff}
\hypersetup{colorlinks=true,linkcolor=SDSoftAccent,urlcolor=SDSoftAccent}
\setlength{\parindent}{0pt}
\setlength{\parskip}{5pt}
\stepdiffsetup{display=notes,theme=soft,layout=lecture,highlight=background,
  reason-style=muted}

\title{\textsf{stepdiff}\\[2pt]\large Visual changes in mathematical derivations}
\author{Ernest Terjyan}
\date{Version 1.1.1 \quad 25 September 2026}

\begin{document}
\maketitle
\vspace{-1.4em}

\begin{abstract}
\noindent
\textsf{stepdiff} is a LuaLaTeX package that highlights visual changes between
successive lines of a mathematical derivation. It is designed for lecture notes,
worked examples, and Beamer slides. It compares printed source tokens; it does
not check whether a mathematical step is valid.
\end{abstract}

\section{First example}

Save the file below as \texttt{example.tex} and compile it with LuaLaTeX.

Put \texttt{stepdiff.sty} and \texttt{stepdiff.lua} beside
\texttt{example.tex}, or install the package in your local TeX tree.

\begin{verbatim}
\documentclass{article}
\usepackage{stepdiff}
\begin{document}
\[
\begin{stepdiff}
\step{a(x+b)}
\step[reason={distribute}]{ax+ab}
\step[reason={commute terms},tag=final]{ab+ax}
\end{stepdiff}
\]
\end{document}
\end{verbatim}

The same source produces this output:
\[
\begin{stepdiff}
\step{a(x+b)}
\step[reason={distribute}]{ax+ab}
\step[reason={commute terms},tag=final]{ab+ax}
\end{stepdiff}
\]
The first line has no predecessor, so automatic highlighting begins on the
second line. The \texttt{reason} appears beside the corresponding step.

\section{A relation-aware example}

The first recognized relation in each line forms an alignment point.
The two sides are compared independently. A changed relation is highlighted
as part of the step.
\[
\begin{stepdiff}[frame=true,color-mode=typed]
\step{a_n \le b_n}
\step[reason={swap sides}]{b_n \ge a_n}
\step[reason={specialize},tag=final]{b_1 \ge a_1}
\end{stepdiff}
\]
The second step swaps both sides, so the comparison is equivalent.
Recognized relation tokens include \(=\), \(\le\), \(\ge\), \(<\), \(>\),
\(\approx\), \(\sim\), \(\equiv\), \(\Rightarrow\), and
\(\Longrightarrow\). Alignment uses the first one found.

\section{Options and controls}

Put global defaults in \verb|\stepdiffsetup{...}| and override them locally:
\begin{verbatim}
\stepdiffsetup{display=notes,theme=soft,highlight=background}
\begin{stepdiff}[frame=true,color-mode=typed]
  ...
\end{stepdiff}
\end{verbatim}

\begin{center}
\small
\begin{tabular}{@{}p{0.25\linewidth}p{0.67\linewidth}@{}}
\textbf{Option} & \textbf{Values and effect}\\
\hline
\texttt{display} & \texttt{notes}, \texttt{slide}, \texttt{focus}: spacing and emphasis presets.\\
\texttt{theme} & \texttt{soft}, \texttt{minimal}, \texttt{focus}: color presets.\\
\texttt{layout} & \texttt{compact}, \texttt{lecture}, \texttt{wide}: column spacing.\\
\texttt{highlight} & \texttt{background}, \texttt{underline}, \texttt{color}, \texttt{none}.\\
\texttt{color-mode} & \texttt{single} (one style), \texttt{typed} (added/modified), \texttt{teaching} (also repeated operations).\\
\texttt{reason-style} & \texttt{plain}, \texttt{muted}, \texttt{badge}.\\
\texttt{frame} & \texttt{true} or \texttt{false}: draw a frame around the derivation.\\
\texttt{legend} & \texttt{true} or \texttt{false}: show typed color categories.\\
\texttt{show-reasons} & \texttt{true} or \texttt{false}: show reason annotations.\\
\end{tabular}
\end{center}

On individual \verb|\step| commands, \texttt{reason=\{...\}} adds an
annotation and \texttt{tag=final} emphasizes a result. The \texttt{diff}
option accepts \texttt{auto} (default), \texttt{none} or \texttt{false}
(suppress automatic highlighting), and \texttt{all} (emphasize the full line).
For example:
\begin{verbatim}
\step[reason={simplify},diff=all,tag=final]{x=2}
\end{verbatim}

\section{Manual correction and visual categories}

Automatic token matching can be approximate for dense notation and complex
macros. To choose exactly what is emphasized, use \texttt{diff=none} and wrap
the desired part of the current line:
\begin{verbatim}
\step[diff=none]{f(x)=\SDchanged{(x+1)^2}}
\end{verbatim}

The public math-safe hooks are \verb|\SDchanged{...}| for a generic change,
\verb|\SDadded{...}| for inserted content, \verb|\SDmodified{...}| for
rewritten content, \verb|\SDoperation{...}| for a conservatively detected
repeated operation, and \verb|\SDfinal{...}| for result emphasis.
\verb|\SDmoved{...}| is available, but move detection is limited.
You can redefine these hooks and \verb|\SDreason{...}| in the preamble;
math hooks must remain safe in math mode.

\section{Slides}

In Beamer, add an overlay specification after \verb|\step|:
\begin{verbatim}
\begin{stepdiff}[display=slide,theme=focus]
\step<1->{f(x)=x^2}
\step<2->[reason={differentiate}]{f'(x)=2x}
\end{stepdiff}
\end{verbatim}
Each step's formula, reason, and highlighting appear together on the
requested overlays. Overlay syntax is accepted in article documents too,
without overlay behavior.

\section{Scope and limitations}

The package compares source tokens visually and does not parse mathematics
semantically or call a computer algebra system. It never validates equality,
inequality, or implication. Removed tokens cannot appear in the new line;
for a deletion-only change, the surviving expression is highlighted as a cue.
If the automatic choice is misleading, use the manual control above.

\section{Customization and troubleshooting}

The generic highlight and reason styles are user hooks. For example, after
loading the package you can change the generic highlight color with:
\begin{verbatim}
\renewcommand{\SDchanged}[1]{%
  \begingroup\color{red!70!black}#1\endgroup}
\end{verbatim}
Keep math styling commands safe in math mode. Other visual hooks are listed
in Section 4; the full set of presentation options is in the README.

If the package is not found, check that both \texttt{stepdiff.sty} and
\texttt{stepdiff.lua} are installed. If an engine error appears, compile
with \texttt{lualatex}. If the first step has no highlight, that is expected:
there is no previous line to compare. For an unexpectedly quiet later step,
check whether \texttt{diff=none} or \texttt{highlight=none} is active.

\section{Support and license}

Source, examples, and issue reporting:
\url{https://github.com/ernestterjyan/stepdiff}.
The package is released under the MIT License; see \texttt{LICENSE}.

\end{document}
